% ===================== ANEXOS =====================
\appendix

\section{C\'odigo Python de C\'alculos}
\label{appendix:codigo}

Los c\'alculos de transferencia de calor y an\'alisis estructural presentados en este informe se implementaron en Python~3.x utilizando las bibliotecas cient\'ificas NumPy y SciPy. El c\'odigo fuente est\'a organizado en los siguientes m\'odulos:

\begin{table}[H]
	\centering
	\caption{M\'odulos de c\'alculo del proyecto P2611.}
	\begin{tabularx}{\textwidth}{@{}lX@{}}
		\toprule
		\textbf{M\'odulo} & \textbf{Descripci\'on} \\
		\midrule
		\texttt{propiedades\_glucosa.py} & Propiedades termof\'isicas de la glucosa Globe~1130 (densidad, viscosidad, $C_p$, $k$) y del agua en funci\'on de la temperatura. \\
		\texttt{geometria\_tanque.py} & Geometr\'ia del tanque: vol\'umenes (cilindro + fondo toriesf\'erico), di\'ametro hidr\'aulico de la media ca\~na, velocidades de flujo. \\
		\texttt{coeficiente\_U.py} & C\'alculo del coeficiente global $U$: correlaciones de Sieder-Tate ($h_i$), Churchill-Chu ($h_o$) y modelo de resistencias en serie. \\
		\texttt{escenarios.py} & Simulaci\'on transitoria del calentamiento (RK45) para los Escenarios~1, 2 y 3. Incluye descarga de 24~ton en el Escenario~1. Generaci\'on de gr\'aficas. \\
		\texttt{ciclo\_descargas.py} & Simulaci\'on del ciclo de 8~descargas a carrotanque (24~h) para los Escenarios~2 y 3. Diagramas de Gantt y tablas por descarga. \\
		\texttt{escenario4\_ciclo.py} & Escenario~4: ciclo industrial de despacho (192~ton, 55~°C, 8~descargas de 24~ton a 57~°C). C\'alculo de tiempos de recalentamiento. \\
		\texttt{espesores\_tanque.py} & An\'alisis estructural: API~650 One-Foot Method (cuerpo) y ASME~VIII UG-32(e) (fondo toriesf\'erico). Cargas de viento y sismo. \\
		\texttt{aislamiento.py} & Dimensionamiento del aislamiento t\'ermico de lana mineral: p\'erdidas de calor y temperatura superficial vs.\ espesor. \\
		\bottomrule
	\end{tabularx}
\end{table}

El c\'odigo fuente completo est\'a disponible en el repositorio del proyecto y puede ser ejecutado de forma independiente para reproducir todos los resultados de este informe.

\section{Propiedades de la Glucosa Globe~1130}
\label{appendix:propiedades}

Las propiedades termof\'isicas de la glucosa Globe~1130 se obtuvieron de las siguientes fuentes:

\begin{itemize}
	\item \textbf{Densidad:} Correlaci\'on lineal ajustada a datos de Perry's Chemical Engineers' Handbook (9th ed.) y Bubble~et~al.~(2007). $\rho(T) = 1430 - 0.50 \cdot T$ [kg/m\textsuperscript{3}], con $T$ en °C.
	\item \textbf{Viscosidad:} Modelo de Arrhenius modificado, ajustado a datos de Telis~et~al.~(2007) para jarabes de carbohidratos de alta concentraci\'on. $\mu(T) = 1.5 \times 10^{-16} \cdot \exp(10\,500 / T_K)$ [Pa$\cdot$s].
	\item \textbf{Calor espec\'ifico:} Correlaci\'on de Choi \& Okos (1986) para soluciones de carbohidratos, ponderada por fracci\'on m\'asica de s\'olidos (83\%).
	\item \textbf{Conductividad t\'ermica:} Correlaci\'on ASHRAE adaptada de Choi \& Okos (1986).
\end{itemize}

\section{Planos de Referencia}
\label{appendix:planos}

Los planos de referencia del tanque (fondo toriesf\'erico y chaqueta de media ca\~na en espiral) se encuentran en la carpeta \texttt{Data/} del proyecto:

\begin{itemize}
	\item \texttt{Data/fondo.png} — Plano del fondo toriesf\'erico con dimensiones.
	\item \texttt{Data/espiral.png} — Plano de la media ca\~na en espiral con paso y perfil.
\end{itemize}

Estos planos fueron proporcionados por INGREDION~SA y constituyen la base dimensional para todos los c\'alculos geom\'etricos de este informe.
